\phantomsection
\pdfbookmark[1]{Аннотация}{Аннотация}

\begin{center}
	\section*{Аннотация}

	выпускной квалификационной работы \\
	студента группы \vkrStudentGroup\space ФГАОУ ВО «МГТУ «СТАНКИН» \\
	\textbf{\vkrStudentR} \\
	по направлению подготовки \\
	\textbf{\vkrMajor «\vkrMajorName»} \\
	на тему: \\
	\textbf{«\vkrTheme»}

\end{center}

\subparagraph{Актуальность работы.} Формирование расписания учебных занятий является одной из основных и наиболее сложных задач автоматизации управления учебным процессом любого учебного заведения. МГТУ «СТАНКИН» не является исключением, каждый семестр специальный отдел в составе Учебно-методического управления МГТУ «СТАНКИН» \cite{stankin:schedule} составляет расписание на следующие полгода, выкладывает его на сайт университета и отправляет преподавателям и старостам групп. На сегодняшний день процесс создания расписания учебных занятий производится вручную, что является медленным и ресурсозатратным процессом. В данной выпускной квалификационной работе будет проведён анализ существующих подходов к решению данной задачи, а также предложен новый алгоритм решения. Финальным шагом работы будет разработка клиент-серверного веб-приложения для создания расписания университета.

\subparagraph{Научная новизна} работы состоит в разработке собственного алгоритма решения задачи составления расписания для учебного заведения ранее не исследованном методом, с возможностью тонкой настройки таких параметров, как процент окон, распределение количества учебных занятий по дням недели и в течение дня и многие другие менее значительные оптимизации.

\subparagraph{Практическая ценность} работы состоит в автоматизации такого ресурсозатратного процесса как планирование и формирование плана учебных занятий на семестр.

\edef\apppages{\fpeval{\getpagerefnumber{vkr:End} - \getpagerefnumber{vkr:AppendixBegin} + 1}}
\edef\totalpages{\getpagerefnumber{vkr:End}}
\edef\figures{\readfromaux{vkrTotalfigures}}
\edef\tables{\readfromaux{vkrTotaltables}}
\edef\books{\readfromaux{vkrTotalbooks}}

\subparagraph{Структура работы.} Выпускная квалификационная работа (ВКР) содержит \totalpages{}  стр.  сквозной нумерации. Количество рисунков в работе --- \figures, таблиц --- \tables. Позиций в списке литературы --- \books.

\subparagraph{Выводы:}

\begin{enumerate}
	\item Проведён анализ методов автоматического составления расписания и обоснован выбор гибридного подхода --- WFC для первичного допустимого решения и эволюционная оптимизация на основе генетического алгоритма для его улучшения.
	\item Разработаны требования к расписанию, описан предложенный алгоритм генерации и оптимизации, спроектирована архитектура и стек технологий клиент-серверного приложения и реализовано программное средство с веб-интерфейсом.
	\item Выполнено экономическое обоснование целесообразности внедрения разработанного средства для МГТУ «СТАНКИН».
\end{enumerate}

\subparagraph{Список работ, опубликованных по теме ВКР:}

\begin{enumerate}
	\item Работа, опубликованная на студенческой научно-практической конференции «Автоматизация и информационные технологии (АИТ-2024)» МГТУ «СТАНКИН» --- «Исследование методов автоматического составления расписания» \cite{ait2024}.
	\item Работа, опубликованная на 68-й Всероссийской научной конференции МФТИ --- «Алгоритм решения задач составления расписания занятий в образовательных учреждениях» \cite{mipt68conf}.
\end{enumerate}
